\section*{Resumen}
\addcontentsline{toc}{section}{Resumen}
\markboth{Resumen}{Resumen}

La adopción de sensores para monitoreo dejó de estar limitada por el costo del equipamiento: un nodo
sensor cuesta entre 5 y 30 dólares y las bandas de radio de uso libre no requieren licencia de
espectro. La barrera se desplazó hacia el destino de los datos. Las organizaciones que ya poseen
sensores descubren que cada dispositivo emplea un formato distinto, que las plataformas comerciales
funcionan únicamente con el hardware de su propio proveedor, y que no disponen de un medio para saber
si sus sensores continúan operando.

Este trabajo propone una plataforma abierta de gestión e interoperabilidad de datos para redes de
sensores heterogéneas, cuyo aporte central es la especificación de un protocolo de comunicación entre
sensor y plataforma basado en el estándar de industria \textbf{Contact ID (ANSI/SIA DC-05)}, empleado
desde hace cuatro décadas por el sector de la seguridad electrónica para supervisar flotas de
dispositivos inalámbricos alimentados a batería. La plataforma extiende su transporte a LoRaWAN,
WiFi/HTTP y MQTT, de modo que organizaciones que ya poseen equipamiento puedan centralizar, normalizar
y administrar sus propios datos.

Esta primera etapa establece la definición del problema mediante un árbol de causas y efectos, formula
un objetivo general y seis objetivos específicos bajo criterios SMART con su matriz de trazabilidad,
cataloga 68 requerimientos funcionales y no funcionales justificados por fuente, evalúa siete
alternativas de protocolo y cinco alternativas de implementación mediante matrices de criterios
ponderados, y construye un backlog de 63 elementos con criterios de aceptación y estimaciones.

La verificación numérica del backlog arrojó un hallazgo determinante para la planificación: el alcance
completo excede en un 37\% el tiempo disponible, mientras que el núcleo demostrativo cabe con un
18\% de holgura. En consecuencia, el compromiso del proyecto se acota explícitamente a ese núcleo y
el resto se declara alcance condicional.

\vspace{0.6em}
\noindent\textbf{Palabras clave:} Internet de las Cosas, interoperabilidad de protocolos, Contact ID,
SIA DC-05, LoRaWAN, gestión de red, PPDIOO.
